\section{Глава 3. Программная реализация системы}


В данной главе представлено описание программной реализации системы и итогового пользовательского интерфейса. По структуре изложения глава построена по модульному принципу: для каждой функциональной подсистемы зафиксировано, что реализовано на клиенте, что реализовано на сервере, как выглядит пользовательский поток и какой результат получает пользователь.

\textbf{[МЕСТО ДЛЯ РИСУНКА 3.1: общий интерфейс приложения (вставляется вручную)]}
Подпись к вставке: набор основных экранов приложения (список поездок, детали поездки, идеи, маршрут, расходы), источник --- актуальная сборка; рисунок должен подтверждать целостность интерфейсного контура.

В качестве демонстрационного контура рассматривается сценарий совместного планирования поездки, где владелец создает поездку, приглашает участника, формирует идеи и маршрут, распределяет расходы, использует погодные данные и AI-подсказки, а также продолжает работу при временной потере сети. Такой выбор позволяет проверить не отдельные экраны, а целостность системы и связность реализованных модулей [7][8][9].

\subsection{3.1 Реализация подсистемы авторизации и профиля}

\subsubsection{Реализация на клиенте}
На стороне Android реализован экран входа и связанный поток инициализации сессии. Пользователь проходит авторизацию и после успешного входа попадает в рабочий контур приложения. В клиентской части также реализован экран настроек, где доступны изменение имени, фотография профиля, завершение сессии и удаление учетной записи в предусмотренном пользовательском сценарии [7][9][14].
В интерфейсе применена политика локализации: при системной локали \texttt{ru} приложение отображает русский UI, в остальных локалях --- английский UI; форматирование дат синхронизировано с той же политикой.

Визуально подсистема построена так, чтобы отделять аутентификационный этап от этапа рабочей навигации. До входа пользователю доступен только поток авторизации; после входа активируется доступ к поездкам и остальным модулям. Такой переход упрощает модель доступа и снижает риск ситуации, когда неавторизованный пользователь может открыть внутренние экраны приложения.

\subsubsection{Реализация на сервере}
На серверной стороне реализованы маршруты аутентификации и сервисы выдачи токенов, включая проверку входных учетных данных и управление жизненным циклом сессии. Сервер отвечает за валидацию входа и разграничение доступа к защищенным маршрутам, связанным с поездками и пользовательскими данными [12][13].

Серверная реализация поддерживает продление/обновление сессии в рамках настроенной политики, что позволяет сохранять предсказуемый пользовательский опыт при длительной работе с приложением. Для операций профиля используются защищенные маршруты, требующие валидного контекста пользователя.

\subsubsection{Пользовательский поток (владелец/участник)}
1. Пользователь выполняет вход в приложении.
2. Система проверяет учетные данные и создает активную сессию.
3. Пользователь открывает настройки и изменяет профильные данные.
4. Обновленные данные сохраняются и отражаются в интерфейсе.

\subsubsection{Результат для пользователя}
Пользователь получает контролируемый доступ к данным поездки, персонализированный профиль и понятный сценарий управления учетной записью. На уровне взаимодействия это снижает риск ошибок доступа и повышает доверие к системе как к личному рабочему пространству.

\textbf{[МЕСТО ДЛЯ РИСУНКА 3.2: авторизация/профиль (вставляется вручную)]}
Подпись к вставке: экран успешного входа и экран настроек профиля после изменения данных; рисунок должен подтверждать выполнение сценария авторизации и управления профилем.

\subsection{3.2 Реализация подсистемы поездок и участников}

\subsubsection{Реализация на клиенте}
Клиентская часть включает экраны списка поездок, формы создания и редактирования поездки, карточку поездки и экран участников. В формах реализованы пользовательские валидаторы дат, названия и базовых параметров поездки, включая валюту и описание. В интерфейсе владельца доступны действия по приглашению участников и управлению составом поездки [7][9].

Отдельное внимание уделено представлению поездки как центра контекста. Пользователь сначала работает со списком поездок, затем открывает детали конкретной поездки и уже из нее переходит в подсистемы идей, маршрута, расходов и уведомлений. Такая навигация уменьшает фрагментацию и делает переходы между задачами последовательными.

\subsubsection{Реализация на сервере}
Сервер реализует маршруты создания, изменения, удаления и выдачи поездок, а также маршруты управления участниками и приглашениями. Для приглашений предусмотрены операции генерации и обработки ссылок, а также проверки валидности приглашения при присоединении. Действия с доступом участников выполняются с учетом роли пользователя в поездке [12][13].

Серверная часть обеспечивает консистентность членства: один и тот же участник не должен быть дублирован в рамках одной поездки, а невалидные приглашения не переводят пользователя в состояние «частично присоединен».

\subsubsection{Пользовательский поток (владелец/участник)}
1. Владелец создает поездку и задает параметры.
2. Владелец формирует приглашение и отправляет ссылку.
3. Участник открывает ссылку и присоединяется к поездке.
4. Оба пользователя видят общую поездку и единый список участников.

\subsubsection{Результат для пользователя}
Группа получает единое рабочее пространство поездки и прозрачный механизм совместного доступа. Это снимает потребность в ручной «сборке команды» вне приложения и сразу переводит участников в общий контекст планирования.

\textbf{[МЕСТО ДЛЯ РИСУНКА 3.3: поездки и участники (вставляется вручную)]}
Подпись к вставке: создание поездки + список участников после присоединения по ссылке; рисунок должен подтверждать сценарий совместного доступа.

\subsection{3.3 Реализация подсистемы идей, обсуждений и маршрута}

\subsubsection{Реализация на клиенте}
На клиенте реализованы экраны списка идей, формы создания/редактирования, карточка идеи с обсуждением, а также экран маршрута по дням. Для обсуждений используется интерфейс, поддерживающий обновление сообщений в реальном времени, что критично для совместного принятия решения. Для маршрута реализованы добавление активностей, изменение порядка, перенос между днями и сценарии обработки дней, выходящих за актуальный диапазон дат поездки [7][9].

Ключевой особенностью интерфейса является связка «идея -> решение -> включение в маршрут». Пользователь не теряет контекст между обсуждением и планом: карточка идеи может быть преобразована в элемент маршрута и дальше участвует в ежедневном планировании поездки.

\subsubsection{Реализация на сервере}
Серверная часть включает маршруты для идей, комментариев и маршрута, а также канал real-time обновлений для обсуждений. Сервер обеспечивает хранение статусов идеи, валидацию операций модерации и согласованную запись изменений в маршрут. В текущей версии real-time канал гарантированно публикует события создания комментариев; удаление комментариев выполняется через отдельный API-маршрут с проверкой прав автора [12][13].

Такое разделение позволяет сохранить единое правило: доменная валидность проверяется на сервере, а клиент отвечает за представление и удобство выполнения действий.

\subsubsection{Пользовательский поток (владелец/участник)}
1. Участник создает идею и публикует ее в поездке.
2. Участники обсуждают идею в комментариях.
3. Владелец принимает решение по идее (одобрение/отклонение).
4. Одобренная идея переносится в конкретный день маршрута.
5. Порядок активностей корректируется в соответствии с планом группы.

\subsubsection{Результат для пользователя}
Пользователь получает непрерывный сценарий от обсуждения к исполнимому маршруту. Это устраняет типичную проблему разрыва между «чатом с предложениями» и «финальным планом», которая характерна для разрозненных инструментов [7][8].

\textbf{[МЕСТО ДЛЯ РИСУНКА 3.4: идеи/обсуждения/маршрут (вставляется вручную)]}
Подпись к вставке: коллаж \texttt{список идей -> обсуждение -> добавление в маршрут}; рисунок должен подтверждать сквозной пользовательский поток.

\subsection{3.4 Реализация подсистемы расходов и балансов}

\subsubsection{Реализация на клиенте}
Клиентская часть включает список расходов, форму создания/редактирования расхода, экран детализации и представление итогового баланса. Поддерживаются сценарии планируемого и оплаченного расхода, выбор плательщика, распределение долей и отметка факта оплаты участниками. Интерфейс ориентирован на то, чтобы пользователь видел не только отдельные записи, но и суммарный финансовый результат по поездке [7][9].

За счет прозрачного отображения долей и статусов оплаты снижается число уточняющих коммуникаций между участниками и упрощается закрытие финансовых обязательств по завершении поездки.

\subsubsection{Реализация на сервере}
Сервер реализует маршруты создания, изменения, удаления и получения расходов. В серверной логике выполняются проверки валидности суммы, участников и связки расхода с конкретной поездкой. Итоговые показатели для отображения баланса вычисляются на клиенте по данным расходов и долей участников, полученных через API [12][13].

\subsubsection{Пользовательский поток (владелец/участник)}
1. Пользователь создает расход с указанием суммы и состава участников.
2. Система рассчитывает доли и обновляет общий финансовый контур.
3. Участники отмечают оплату в рамках своих обязательств.
4. Баланс пересчитывается и отражается в интерфейсе.

\subsubsection{Результат для пользователя}
Группа получает прозрачную и проверяемую финансовую модель поездки. Пользователь видит, кто оплатил, кто должен и как меняется итоговый баланс при каждом действии.

\textbf{[МЕСТО ДЛЯ РИСУНКА 3.5: расходы и баланс (вставляется вручную)]}
Подпись к вставке: экран расходов с итоговыми суммами и деталями долей; рисунок должен подтверждать пересчет баланса после изменения расхода.

\subsection{3.5 Реализация подсистем погоды, AI-подсказок и уведомлений}

\subsubsection{Реализация на клиенте}
На клиенте реализованы экран прогноза и погодные блоки в контексте поездки, интерфейс генерации AI-подсказок и экран настроек уведомлений. Пользователь может обновлять погодные данные, выбирать релевантный контекст поездки для AI-предложений и сохранять подходящие варианты в список идей. В настройках доступны переключатели категорий уведомлений [7][9].

Такая интеграция делает вспомогательные сервисы частью рабочего процесса, а не внешним шагом «после планирования».

\subsubsection{Реализация на сервере}
Серверная часть интегрируется с погодным провайдером и AI-провайдером через выделенные маршруты и сервисные адаптеры. Для уведомлений реализованы маршруты регистрации токенов устройств и отправки событий. При этом погодные и AI-данные не заменяют базовые доменные функции поездки, а дополняют их контекстом принятия решений [17][18][20].

\subsubsection{Пользовательский поток (владелец/участник)}
1. Пользователь открывает погодный модуль поездки и обновляет прогноз.
2. Пользователь запрашивает AI-предложения для маршрута.
3. Подходящее предложение сохраняется в идеи поездки.
4. При включенных категориях уведомлений, выданном системном разрешении и доступном интернете пользователь получает уведомления по событиям идей (создание/комментарии) и расходам.

\subsubsection{Результат для пользователя}
Пользователь быстрее принимает решения по маршруту и в поддерживаемых условиях доставки push видит актуальные события поездки без постоянной ручной проверки каждого экрана.

\subsection{3.6 Реализация офлайн-режима и синхронизации}

\subsubsection{Реализация на клиенте}
В клиентской части реализованы локальное хранение и очередь изменений, позволяющие выполнять операции создания, изменения и удаления поддерживаемых сущностей при отсутствии сети. Изменения, попавшие в очередь, помечаются для последующей отправки после восстановления подключения [9][19].

Синхронизация запускается фоновыми механизмами и при стандартной активности в приложении. Для пользователя этот процесс выглядит как восстановление актуального состояния без повторного ручного ввода данных в типовых сценариях.

\subsubsection{Реализация на сервере}
Сервер принимает пакет изменений синхронизации для операций `create/update/delete`, обрабатывает элементы изолированно и применяет серверную стратегию разрешения коллизий, формируя согласованное итоговое состояние данных. Такой подход позволяет завершать частично корректный пакет без отклонения всего запроса и без ручного разрешения конфликтов на стороне пользователя [12][13][19].

\subsubsection{Пользовательский поток (владелец/участник)}
1. Пользователь создает, изменяет или удаляет данные без сети.
2. Изменения сохраняются локально в очередь синхронизации.
3. После восстановления сети изменения отправляются на сервер.
4. Интерфейс обновляется и показывает единое актуальное состояние.

\subsubsection{Результат для пользователя}
Система поддерживает непрерывность работы для офлайн-операций создания, изменения и удаления, а также снижает вероятность потери изменений в мобильных условиях нестабильного подключения за счет серверного разрешения коллизий при синхронизации.

\textbf{[МЕСТО ДЛЯ РИСУНКА 3.6: офлайн и синхронизация (вставляется вручную)]}
Подпись к вставке: коллаж \texttt{действие без сети} и \texttt{результат после восстановления сети}; рисунок должен подтверждать применение офлайн-очереди.

\subsection{3.7 Проверка сквозных пользовательских сценариев}

В рамках верификации результата проведена функциональная проверка сквозных сценариев, отражающих ядро пользовательской ценности приложения. Проверка ориентирована на связность модулей и достижение пользовательского результата, а не только на изолированную работоспособность отдельных экранов.

Текущее покрытие модулей по реализованным возможностям представлено в таблице.

\begin{table}[H]
\centering
\small
\caption{Покрытие модулей приложения}
\begin{tabular}{|p{2.8cm}|p{5.4cm}|p{2.8cm}|p{1.9cm}|}
\hline
Модуль & Реализованные возможности & Клиент/Сервер & Статус \\ 
\hline
Авторизация и профиль & Вход, изменение профиля, завершение сессии, удаление аккаунта в пользовательском потоке & Android + Backend & Реализовано \\ 
\hline
Поездки и участники & Создание/редактирование/удаление поездки, приглашение и присоединение & Android + Backend & Реализовано \\ 
\hline
Идеи и обсуждения & CRUD идей, статусы, комментарии в реальном времени & Android + Backend + WS & Реализовано \\ 
\hline
Маршрут & Планирование по дням, перенос и переупорядочивание активностей & Android + Backend & Реализовано \\ 
\hline
Расходы и баланс & Ведение расходов, сплиты, баланс и статусы оплаты & Android + Backend & Реализовано \\ 
\hline
Погода, AI, уведомления & Прогноз, генерация AI-предложений, настройки уведомлений & Android + Backend + внешние сервисы & Реализовано \\ 
\hline
Офлайн и синхронизация & Локальная очередь изменений и фоновая синхронизация & Android + Backend & Реализовано \\ 
\hline
\end{tabular}
\end{table}

Результаты проверки сквозных сценариев приведены в таблице ниже.

\begin{table}[H]
\centering
\small
\caption{Проверка сквозных пользовательских сценариев}
\begin{tabular}{|p{3.2cm}|p{3.3cm}|p{3.8cm}|p{2.6cm}|}
\hline
Сценарий проверки & Шаги & Ожидаемый результат & Фактический результат \\ 
\hline
Создание поездки и приглашение участника & Владелец создает поездку, отправляет приглашение, участник присоединяется & У обоих пользователей отображается одна и та же поездка и общий состав участников & Сценарий воспроизводится на контрольном наборе, критических отклонений не зафиксировано \\ 
\hline
Идея $\rightarrow$ обсуждение $\rightarrow$ маршрут & Участник создает идею, обсуждение проходит в реальном времени, идея переносится в маршрут & Идея из обсуждения становится элементом маршрута выбранного дня & Сценарий воспроизводится; результат подтвержден в пределах проверенного набора \\ 
\hline
Управление расходом и балансом & Создается расход, задаются доли, фиксируется оплата & Баланс пересчитывается на клиенте по данным расходов и долей участников & Сценарий воспроизводится; несоответствий в проверенном наборе не обнаружено \\ 
\hline
Погода + AI-подсказка & Пользователь обновляет прогноз, генерирует AI-предложение, сохраняет в идеи & Предложение появляется в идеях, погодные данные отображаются по контексту поездки & Сценарий воспроизводится при доступности внешних сервисов \\ 
\hline
Офлайн-изменение и синхронизация & Действия выполняются без сети, затем сеть восстанавливается & Офлайн-очередь для операций создания/изменения/удаления отправляется после восстановления сети; коллизии разрешаются сервером & Сценарий воспроизводится на контрольном наборе; итоговое состояние согласуется после синхронизации \\ 
\hline
Настройки уведомлений & Пользователь изменяет категории уведомлений в настройках & При включенных категориях, выданном системном разрешении и доступном интернете система применяет выбранные категории (события идей и расходов) & Сценарий воспроизводится; применяются только поддерживаемые категории \\ 
\hline
\end{tabular}
\end{table}

\textbf{[МЕСТО ДЛЯ РИСУНКА 3.7: пример проверки сценария (вставляется вручную)]}
Подпись к вставке: фрагмент протокола ручного сценария (шаги + факт выполнения); рисунок должен подтверждать воспроизводимость контрольной проверки.

Проверка выполнялась как ручной сценарный прогон в отладочном контуре приложения с использованием контрольного набора пользовательских ситуаций. Представленные результаты отражают покрытие ключевых сценариев в рамках этого набора и не претендуют на исчерпывающую верификацию всех возможных комбинаций входных данных [7][8][9].

Проведенная проверка показывает, что реализованная система закрывает целевой набор пользовательских сценариев, сформулированный в ТЗ и обоснованный в аналитической части работы [7][8]. При этом дальнейшим этапом развития остается расширение автоматизированного регрессионного покрытия для сокращения стоимости последующих изменений.

\subsection{3.8 Автоматизированное тестирование и анализ покрытия}

Метрики раздела 3.8 используются как инженерный контроль процесса разработки (Android + backend) и не заменяют клиентский приемочный контур ПМИ, который проверяет только пользовательские сценарии взаимодействия с приложением.

Для снижения регрессионных рисков в проекте внедрены автоматизированные модульные тесты и JVM-тесты пользовательского интерфейса (Compose/Robolectric), а также единые пороговые критерии качества по покрытию в клиентской и серверной частях [9].

\subsubsection{Пошаговый запуск проверок (для внешнего проверяющего)}

Чтобы воспроизвести результаты без знания внутренней структуры проекта, используется следующий порядок:
1. открыть корневую директорию репозитория, где находятся каталоги \texttt{android/}, \texttt{backend/} и скрипт \texttt{quality-check.sh};
2. запустить единый сценарий проверки:

\begin{verbatim}
./quality-check.sh
\end{verbatim}

3. при необходимости запускать модули отдельно:

\begin{verbatim}
# серверный модуль
cd backend && ./gradlew qualityCheck

# клиентский модуль Android
cd android && ./gradlew :app:qualityCheck
\end{verbatim}

Сценарий \texttt{qualityCheck} в каждом модуле включает:
1. выполнение тестов;
2. формирование отчета JaCoCo в форматах \texttt{XML} и \texttt{HTML};
3. проверку порогов покрытия строк и покрытия ветвлений.

Артефакты проверки находятся в фиксированных путях:
1. серверный модуль:
\texttt{backend/build/test-results/test/*.xml},
\texttt{backend/build/reports/jacoco/test/jacocoTestReport.xml},
\texttt{backend/build/reports/jacoco/test/html/index.html};
2. клиентский модуль Android:
\texttt{android/app/build/test-results/testDebugUnitTest/*.xml},
\texttt{android/app/build/reports/jacoco/jacocoDebugUnitTestReport/jacocoDebugUnitTestReport.xml},
\texttt{android/app/build/reports/jacoco/jacocoDebugUnitTestReport/html/index.html}.

Критерий прохождения: в каждом модуле одновременно выполняются пороги \texttt{покрытие строк >= 50\%} и \texttt{покрытие ветвлений >= 40\%}. При нарушении порога задача Gradle завершается ошибкой, что обеспечивает применимость проверки в контуре CI.

\subsubsection{Методика расчета коэффициента покрытия}

В работе используется JaCoCo. Для каждого счетчика рассчитывается:

\texttt{КоэффициентПокрытия = Покрыто / (Покрыто + Непокрыто) * 100\%}

Ключевые метрики:
1. \texttt{ПокрытиеСтрок = ПокрытыеСтроки / (ПокрытыеСтроки + НепокрытыеСтроки) * 100\%}
2. \texttt{ПокрытиеВетвлений = ПокрытыеВетви / (ПокрытыеВетви + НепокрытыеВетви) * 100\%}

Где:
1. \texttt{ПокрытыеСтроки} --- строки, для которых в ходе теста выполнена хотя бы одна инструкция;
2. \texttt{НепокрытыеСтроки} --- строки, для которых не выполнено ни одной инструкции;
3. \texttt{ПокрытыеВетви/НепокрытыеВетви} --- покрытые/непокрытые ветви условных конструкций.

\subsubsection{Контур файлов и причины исключений из расчета покрытия}

Покрытие считается не по всему бинарному артефакту, а по целевому доменному контуру, заданному в Gradle-конфигурации [9]. Такой подход нужен, чтобы метрика отражала качество тестирования бизнес-логики, а не технического каркаса.

Для клиентского модуля Android в расчет включаются скомпилированные классы \texttt{nvk/cotrip/**} (сборка \texttt{debug}), при этом исключаются:
1. сгенерированные и инфраструктурные классы (\texttt{R}, \texttt{BuildConfig}, \texttt{Manifest}, DI/Hilt/Dagger, \texttt{*\_Factory}, \texttt{*\_MembersInjector}, \texttt{*\_Impl});
2. обвязка пользовательского интерфейса и точки входа (\texttt{*ScreenKt*}, \texttt{MainActivity}, \texttt{CoTripApp});
3. отдельные инфраструктурные директории (\texttt{di/**}, \texttt{ui/theme/**}, \texttt{data/network/dto/**});
4. тестовые и синтетические классы (\texttt{*Test*}, часть \texttt{*\$*} шаблонов).

Для серверного модуля покрытие считается по классам основного набора (\texttt{main}) с исключением инфраструктурных слоев:
1. точки входа и загрузочная инициализация (\texttt{ApplicationKt}, часть подсистемы журналирования и конфигурации);
2. маршруты и транспортный слой (\texttt{routes/**}, \texttt{plugins/**});
3. инфраструктура БД (\texttt{db/**}, \texttt{DatabaseFactory});
4. часть внешних адаптеров доставки событий (\texttt{FirebasePushService}, \texttt{CommentEventsPublisher}, \texttt{CommentsHub});
5. служебные/сгенерированные вспомогательные конструкции Kotlin (\texttt{\$Companion}, \texttt{\$serializer}, \texttt{\$WhenMappings}, \texttt{\$inlined\$}).

Причины исключения указанных групп кода:
1. сгенерированный код не содержит проектных решений и может существенно «шуметь» в метрике, искажая сравнение между итерациями;
2. код начальной инициализации и внедрения зависимостей в основном выполняет связывание компонентов, а не реализует прикладные правила предметной области;
3. транспортный слой и внешние адаптеры проверяются преимущественно интеграционными и сценарными проверками, где приоритет имеет корректность контракта в целом, а не покрытие строк отдельных оберток;
4. исключение синтетических конструкций Kotlin предотвращает техническое занижение покрытия на коде, который разработчик напрямую не пишет.

Исключение из расчета покрытия не означает отсутствия проверки: часть таких компонентов контролируется косвенно через интеграционные тесты и сквозные пользовательские сценарии (раздел 3.7). Метрика покрытия в данном контуре используется как индикатор качества автоматизированных тестов бизнес-логики.

\subsubsection{Фактические результаты}

Количественные показатели получены по результатам запуска \texttt{./quality-check.sh} от 18.03.2026 в текущем состоянии репозитория.

\begin{table}[H]
\centering
\small
\caption{Результаты автоматизированного тестирования и покрытия}
\begin{tabular}{|p{2.0cm}|p{1.0cm}|p{1.2cm}|p{1.0cm}|p{2.1cm}|p{2.1cm}|p{1.6cm}|}
\hline
Модуль & XML-отчеты & Тесты & Пропуски & Покрытие строк & Покрытие ветвлений & Пороговый контроль \\ 
\hline
Клиент Android & 61 & 291 (0 ошибок) & 0 & 3559/5677 (62.69\%) & 1597/3918 (40.76\%) & Пройден \\ 
\hline
Сервер & 10 & 46 (0 ошибок) & 0 & 534/968 (55.17\%) & 280/577 (48.53\%) & Пройден \\ 
\hline
Итог (сумма) & 71 & 337 (0 ошибок) & 0 & 4093/6645 (61.60\%) & 1877/4495 (41.76\%) & Пройден \\ 
\hline
\end{tabular}
\end{table}

Пороговые значения в обоих модулях: покрытие строк не ниже 50\%, покрытие ветвлений не ниже 40\% (оба условия выполнены).

Дополнительно по контурным объектам JaCoCo:
1. клиент Android: 32 пакета, 223 исходных файла, 192 класса в расчетном контуре (153 покрыто полностью/частично, 39 не покрыто).
2. сервер: 7 пакетов, 14 исходных файлов, 66 классов в расчетном контуре (51 покрыто полностью/частично, 15 не покрыто).

\subsubsection{Интерпретация результатов}

Полученные значения подтверждают выполнение минимальных пороговых критериев качества в обоих модулях. При этом:
1. покрытие строк в обоих модулях превышает установленный порог и отражает приемлемый охват исполняемых сценариев;
2. покрытие ветвлений в Android-модуле находится близко к порогу, что указывает на потенциал расширения набора тестов по альтернативным ветвям условий и граничным состояниям;
3. по зафиксированному прогону пропуски отсутствуют, что упрощает интерпретацию регрессионных метрик при повторных запусках.

Таким образом, автоматизированная проверка дополняет ручную верификацию сквозных сценариев (раздел 3.7) и формирует воспроизводимый контур контроля качества для последующих итераций разработки [9].

\subsection{3.9 Выводы по главе}

В третьей главе представлена программная реализация приложения в модульной форме, соответствующей структуре и стилю рассмотренных примеров ВКР. Для каждой подсистемы зафиксированы реализованные возможности на клиенте и сервере, а также итоговый пользовательский результат.

Показано, что ключевая ценность системы формируется на уровне сквозных сценариев: участники группы могут совместно пройти путь от создания поездки и обсуждения идей до построения маршрута, распределения расходов и работы в условиях нестабильной сети. Реализация вспомогательных модулей (погода, AI, уведомления) встроена в общий процесс планирования и усиливает его практическую применимость.

Результаты функциональной проверки на контрольном наборе сценариев указывают на соответствие текущей версии приложения целевому контуру совместного планирования, заявленному в предыдущих главах. Таким образом, после аналитического и проектного этапов в главах 1 и 2, в главе 3 получено практическое подтверждение применимости разработанной системы в пределах выполненной проверки.
